% 本文件由 LaTeX 排版 Agent 根据有效 translation.json 自动生成。
% author=Papias; work_id=0125; title=帕皮亚残篇

\chapter{\Person{帕皮亚}残篇}

% structure: p0001 source=h1 target=omitted reason=page-title-already-rendered-by-parent

% structure: p0002 source=h2 target=section reason=relative-heading-rank; base=h2

\section{一}

出自\WorkTitle{主圣言的阐释}。\par

\EditorialNote{帕皮亚的作品在通常流传中计有五卷，并被称为\WorkTitle{主圣言的阐释}。依勒内提及这些是他唯一著作，如下文所述：\Quote{现在这些事由\Person{帕皮亚}——一位古代之人，曾听\Person{若望}讲道，并与\Person{波利卡普}为友——在他的第四卷书中书面见证；因为他共写了五卷书。}\Person{依勒内}如此写道。此外，\Person{帕皮亚}本人在其著作的引言中表明，他本人并非圣宗徒的听者与目击者；但他告诉我们，他从那些熟悉宗徒的人那里接受了我们宗教的真理，如下文所述：}\par

但我并不介意在我的阐释之外，也将我从\Person{长老们}那里随时谨慎接受的训诲，并小心记忆于脑海中的，同时向你们保证其真实。因为我并不像众人那样，喜欢说话多的人，而是喜欢教导真理的人；也不是喜欢讲述奇异诫命的人，而是喜欢复述\Person{主}为信德所赐、并由真理本身所出的诫命。因此，若有任何曾跟随\Person{长老}的人来，我仔细询问他们的说法——\Person{安德肋}或\Person{伯多禄}说了什么，或是\Person{斐理伯}、\Person{多默}、\Person{雅各伯}、\Person{若望}、\Person{玛窦}，或\Person{主}的任何其他门徒说了什么；以及\Person{主}的门徒\Person{阿里斯蒂昂}和\Person{长老若望}所说的事。因为我想，从书本获得的东西，不如从活生生而常存的声音中获得的那样对我有益。\par

% structure: p0006 source=h2 target=section reason=relative-heading-rank; base=h2

\section{二}

\EditorialNote{初期基督徒}\Emph{称那些按神意履行淳朴无伪的人}为\Emph{孩子}，\EditorialNote{正如\Person{帕皮亚}在\WorkTitle{主圣言的阐释}第一卷中所说，以及\Person{克莱孟·亚历山大}在他的\WorkTitle{导师}中所说。}\par

% structure: p0008 source=h2 target=section reason=relative-heading-rank; base=h2

\section{三}

\Person{犹达斯}在这世界上行走，成为不虔的悲惨榜样；因为他的身体肿胀到无法通过战车能轻易通过的地方，他被战车压碎，以致他的肠子迸流而出。\par

% structure: p0010 source=h2 target=section reason=relative-heading-rank; base=h2

\section{四}

\EditorialNote{正如曾看见\Person{主}的门徒\Person{若望}的长老们所记，他们曾从他那里听到\Person{主}关于那些时代如何教导，并说：}\Quote{日子将到，葡萄树要生长，每棵有一万根枝条，每根枝条有一万根细枝，每根真细枝有一万条嫩梢，每条嫩梢上有一万个葡萄串，每个葡萄串上有一万粒葡萄，每粒葡萄压榨后可得二十五\OriginalTerm{默勒特斯}{metretes}酒。当任何一位圣人抓住一个葡萄串时，另一个将喊道：\InnerQuote{我是更好的葡萄串，拿我；藉着我赞美\Person{主}。}同样，他说，一颗麦粒将生产一万个麦穗，每个麦穗有一万粒麦子，每粒麦子产出十磅清纯精面粉；苹果、种子和草也将以类似的比例产出；所有动物，那时只吃地上的出产，将变得和平和谐，完全服从于人。}\EditorialNote{这些事由\Person{帕皮亚}——一位古代之人，曾听\Person{若望}讲道，并与\Person{波利卡普}为友——在他的第四卷书中书面见证；因为他共写了五卷书。他又补充说：\Quote{这些事对信徒是可信的。而叛徒\Person{犹达斯}，}他说，\Quote{不信，并问：\InnerQuote{\Person{主}怎能成就这样的生长？}\Person{主}说：\InnerQuote{那些来到它们那里的人将会看见。}}这些就是先知\Person{依撒意亚}所说的时候：\InnerQuote{豺狼将与羔羊同卧}等等。\BibleRef{依 11:6}}\par

% structure: p0012 source=h2 target=section reason=relative-heading-rank; base=h2

\section{五}

正如长老们所说，那么那些被认为配得上天堂居所的人将去那里，另一些人将享受乐园的愉悦，还有一些人将拥有城市的荣光；因为\Person{救主}将在各处显现，依照看见他的人所配得的。但是，那些结百倍果实者、六十倍果实者和三十倍果实者的居所有区别；因为第一类将被提升到天上，第二类将居住在乐园，最后一类将住在城市；而因此\Person{主}说：\Quote{在我\Person{父}的家里有许多住处：}\BibleRef{若 14:2}因为万有都属于\Person{天主}，祂供给万物合适的居所，正如祂的话语所说，每个人都由\Person{父}分得一份，按照每个人是或将是怎样的人。而这就是\BibleRef{玛 22:10}坐席，受邀请参加婚宴的人将躺卧其上。宗徒的门徒——长老们说，这就是被救者的等级和安排，并且他们通过这样的步骤前进；而且，他们通过\Person{圣神}上升到\Person{圣子}，并通过\Person{圣子}上升到\Person{圣父}；并且到适当的时候，\Person{圣子}将把祂的工作交给\Person{圣父}，正如\Person{宗徒}所说：\Quote{因为祂必须为王，直到把一切仇敌放在祂脚下。最后被毁灭的仇敌是死亡。}\BibleRef{格前 15:25}因为在王国时期，在地上的义人将忘记死亡。\Quote{但祂说万物都放在祂下，显然，那将万物放在祂下的除外。及至万物都顺服了祂，那时\Person{子}自己也要顺服那叫万物顺服祂的，好叫\Person{天主}成为万物中的万有。}\BibleRef{格前 15:27}\par

% structure: p0014 source=h2 target=section reason=relative-heading-rank; base=h2

\section{六}

［\Person{Papias}（即我们刚才提及的那位）宣称，他从那些陪伴宗徒的人那里领受了宗徒的言论；他还声称，他曾亲耳听到\Person{Aristion}和长老\Person{约翰}的教导。因此他屡次在著作中提及他们，并记下他们的传承。我们提及这些情况或许不无益处。在已引述的\Person{Papias}的陈述之外，再补充他的一些其他段落，其中记载了一些奇异的事迹，并说明他是从传统中获得这些知识的。关于宗徒\Person{斐理伯}与其女儿们住在\Place{希拉波利斯}的事迹，上文已经提到。现在我们应当指出，同时代的\Person{Papias}是如何记载他从\Person{斐理伯}的女儿们那里听到的奇妙叙述的。因为他记载说，在他那个时代，一个死人曾复活。他还提到另一个关于\Person{犹斯托}（别号\Person{巴尔撒巴}）的奇迹：他因主的恩宠，喝下致命的毒药而毫发无损。此外，同一个人还记录了一些来自未成文的传统的事，其中包括救主的一些奇异的比喻和训诲，以及一些更具传奇性质的内容。在这些内容中，他说，从死者中复活之后，将有一个千年王国，那时基督将亲自在地上为王。他还在他自己的作品中，记录了前面提到的\Person{Aristion}所传的主的言论，以及长老\Person{约翰}的传承。关于这些，我们只能请读者自行参阅原书；但如今，在已有的摘录之外，我们将添加一项最重要的事，即他［\Person{Papias}］关于撰写福音的\Person{马尔谷}的传统，其记载如下］：长老这样说：\Person{马尔谷}既成为\Person{伯多禄}的译员，便准确地记下了他所记得的一切。然而，他并未按顺序记述基督的言行。因为他既没有听过主，也没有伴随过主。但后来，如我所说，他伴随了\Person{伯多禄}，而\Person{伯多禄}的教导是为了适应听众的需要，并无意系统地记述主的言论。因此，\Person{马尔谷}这样凭记忆写下一些事情，并没有错。因为他特别留意一件事：不遗漏他所听到的任何内容，也不在陈述中掺入虚假。［这便是\Person{Papias}关于\Person{马尔谷}的记载；而关于\Person{玛窦}，他作了如下陈述］：\Person{玛窦}用希伯来语收集了［主的］神谕，而每个人都尽己所能地解释它们。［同一个人还引用\BibleBook{若望}{一书}和\BibleBook{伯多禄}{前书}中的证据。此外，他还讲述了另一个关于一个妇人在主面前被指控许多罪的故事，这故事见于\WorkTitle{希伯来人福音}。］\par

% structure: p0016 source=h2 target=section reason=relative-heading-rank; base=h2

\section{VII}

\Emph{\Person{Papias}这样逐字说道：}\Quote{祂将掌管世界秩序的权力赐给了他们中的一些［天使］，并委派他们善用其权力。}\Emph{紧接着他又说：}\Quote{但他们的秩序归于无有了。}\par

% structure: p0018 source=h2 target=section reason=relative-heading-rank; base=h2

\section{VIII}

关于这本书（\WorkTitle{默示录}）的灵感，我们认为无需赘言；因为有福的\Person{额我略·神学家}和\Person{济利禄}，甚至更早的\Person{Papias}、\Person{依勒内}、\Person{梅多迪乌斯}和\Person{希玻里}，都已经为此作了完全充分的见证。\par

% structure: p0020 source=h2 target=section reason=relative-heading-rank; base=h2

\section{IX}

藉\Place{希拉波利斯}著名的\Person{Papias}（他是那位曾靠在基督胸怀的宗徒的门徒）、\Person{克莱孟}、亚历山大里亚［教会］的司铎\Person{庞代努斯}，以及智慧的\Person{阿摩尼乌斯}这些古代的最初注释者，他们彼此意见一致，将六天的创造工程理解为指向基督和整个教会。\par

% structure: p0022 source=h2 target=section reason=relative-heading-rank; base=h2

\section{X}

(1.) 主之母\Person{玛利亚}；(2.) \Person{克洛帕}或\Person{阿耳斐}的妻子\Person{玛利亚}，她是主教兼宗徒的\Person{雅各伯}、\Person{西满}、\Person{达陡}和一位\Person{若瑟}的母亲；(3.) \Person{载伯德}的妻子\Person{玛利亚·撒罗默}，是福音作者\Person{若望}和\Person{雅各伯}的母亲；(4.) \Person{玛利亚·玛达肋纳}。这四位见于福音书中。\Person{雅各伯}、\Person{犹达}和\Person{若瑟}是主的一位姑母（2）的儿子。\Person{雅各伯}和\Person{若望}是主的另一位姑母（3）的儿子。\Person{玛利亚}（2）——小\Person{雅各伯}和\Person{若瑟}的母亲、\Person{阿耳斐}的妻子——是主之母\Person{玛利亚}的姐妹，\Person{若望}称她为\Person{克洛帕}的\Person{玛利亚}，或因她的父亲，或因家族宗族，或其他原因。\Person{玛利亚·撒罗默}（3）被称为\Person{撒罗默}，或因其丈夫，或因其村庄。有些人认为她与\Person{克洛帕}的\Person{玛利亚}是同一人，因为她曾有两个丈夫。\par

\SourceRecord{Translated by Alexander Roberts and James Donaldson.；Ante-Nicene Fathers；Vol. 1.；Edited by Alexander Roberts, James Donaldson, and A. Cleveland Coxe.；Buffalo, NY: Christian Literature Publishing Co.,；1885.；<http://www.newadvent.org/fathers/0125.htm>.}
